\chapter{理论基础与研究现状}
\section{人体几何重建理论基础}
\vspace{0.3\baselineskip}
\subsection{三维人体模型的表示方式}
三维人体模型具有多种标准表示方式，常见的表示方式包括网格、点云、体素、深度图等显式表示方式（见图2­2）和占用场、符号距离场等隐式表示方式（见图2­3），不同的表示方式在几何精度、计算效率和应用适配性等方面的表现不同。网格是当前三维人体模型应用的最为直接的表示方式，在模型驱动和动画创作等领域有着广泛应用。然而，由于网格结构中顶点与面片的数据分布呈非规则形态，使得其难以被深度学习模型直接建模与预测。因此，现有的三维人体重建方法往往引入其他类型的中间表示，作为过渡手段，以间接完成对网格数据的建构与恢复。图2­1展示了不同表示方法之间相互转换的关系。本节主要阐述最常用的三维数据表示方法及其优缺点，并分析其在人体三维重建任务的适配程度。
这是正文的第一章第一节。根据规范，正文采用小四号宋体，行距为1.5倍。
这里测试一下英文和数字的字体：The standard requires Times New Roman for English and numbers, e.g., $E=mc^2$ or 2025.

\subsection{三维人体模型的表示方式}
123这是正文的第一章第一节。根据规范，正文采用小四号宋体，行距为1.5倍。
这里测试一下英文和数字的字体：The standard requires Times New Roman for English and numbers, e.g., $E=mc^2$ or 2025.

\subsection{参数化人体模型}
123这是正文的第一章第一节。根据规范，正文采用小四号宋体，行距为1.5倍。
123这是正文的第一章第一节。根据规范，正文采用小四号宋体，行距为1.5倍。
这里测试一下英文和数字的字体：The standard requires Times New Roman for English and numbers, e.g., $E=mc^2$ or 2025.
这里测试一下英文和数字的字体：The standard requires Times New Roman for English and numbers, e.g., $E=mc^2$ or 2025.123这是正文的第一章第一节。根据规范，正文采用小四号宋体，行距为1.5倍。
这里测试一下英文和数字的字体：The standard requires Times New Roman for English and numbers, e.g., $E=mc^2$ or 2025.
\subsection{人体网格提取算法}
123这是正文的第一章第一节。根据规范，正文采用小四号宋体，行距为1.5倍。
这里测试一下英文和数字的字体：The standard requires Times New Roman for English and numbers, e.g., $E=mc^2$ or 2025.123这是正文的第一章第一节。根据规范，正文采用小四号宋体，行距为1.5倍。

这里测试一下英文和数字的字体：The standard requires Times New Roman for English and numbers, e.g., $E=mc^2$ or 2025.

\section{人体几何重建研究现状}
这是二级标题（小三号黑体）。

\subsection{基于参数化模型的方法}
123
\subsection{基于多识图几何的方法}
123
\subsection{基于模版的方法}
123
\subsection{基于尹是函数的方法}
123

\section{本章小结}
这是三级标题（四号黑体）。
当文字填满一页时，你可以观察页眉在奇数页和偶数页的变化。
